\section*{Chapitre \ref{ch:b3:lab-techniques-3} --- Techniques de laboratoire III~: la pratique de la recherche}

\begin{solution}{exo:b3:lab-techniques-3:1}
$(0.50/1013)^3 = \num{1.2e-10}$. Les fuites, le gaz adsorbé sur le verre et le dégazage de la graisse et des septums
en laissent bien davantage~; une verrerie sèche et chaude et de bons rodages
comptent plus que des cycles supplémentaires.
\end{solution}

\begin{solution}{exo:b3:lab-techniques-3:2}
$\qty{20.0}{mmol}/\qty{1.48}{mol/L} = \qty{13.5}{mL}$.
\end{solution}

\begin{solution}{exo:b3:lab-techniques-3:3}
$120 + 10 \times 1.00782503 + 55.93493633 - 0.00054858 = \qty{186.01264}{u}$. Erreur~: $10^6 \times (186.0129 - 186.01264)/186.01264 = +1.4$ ppm,
bien en deçà de 5 ppm.
\end{solution}

\begin{solution}{exo:b3:lab-techniques-3:4}
\omterm{def:b3:lab-techniques-3:literature}{Littérature primaire}~: l'article de périodique, le brevet, la thèse.
Secondaire~: l'article de synthèse, le recueil de constantes. Tertiaire~:
le manuel.
\end{solution}

\begin{solution}{exo:b3:lab-techniques-3:5}
$M = \qty{194.0}{g/mol}$~: C $96.0/194.0 = 49.48$\,\%, H $10.0/194.0 = 5.15$\,\%, N $56.0/194.0 = 28.87$\,\%. Écarts de 0.17, 0.07 et 0.16
point~: à moins de 0.4, l'analyse confirme la formule.
\end{solution}

\begin{solution}{exo:b3:lab-techniques-3:6}
$\qty{170.0}{mg}/\qty{212.0}{g/mol} = \qty{0.802}{mmol}$~; $c = \qty{0.802}{mmol}/\qty{0.54}{mL} = \qty{1.5}{mol/L}$ (deux chiffres significatifs, imposés par
le volume lu).
\end{solution}

\begin{solution}{exo:b3:lab-techniques-3:7}
Les moyennes valent 73.0 et 77.0\,\%~; les deux écarts types valent $\sqrt{10/4} = 1.58$~; $s_p = 1.58$. $t = 4.0/(1.58\sqrt{2/5}) = 4.0$, supérieur à
2.31~: l'amélioration est significative au niveau de 95\,\%.
\end{solution}

\begin{solution}{exo:b3:lab-techniques-3:8}
$F = (0.80/0.30)^2 = 7.1 > 5.05$~: la première méthode est significativement moins fidèle.
\end{solution}

\begin{solution}{exo:b3:lab-techniques-3:9}
Acheter de petits flacons avec inhibiteur, dater chacun à la réception et
à l'ouverture, les stocker fermés, à l'abri de la lumière et au frais~;
tester les peroxydes avant toute distillation ou évaporation et à
intervalles fixes après ouverture~; éliminer les flacons périmés, ou
positifs, par la filière des déchets dangereux sans les concentrer.
\end{solution}

\begin{solution}{exo:b3:lab-techniques-3:10}
Dangers~: dégagement de dihydrogène (incendie, pression), réactivité de
NaH avec l'eau, et décomposition exothermique connue du DMF avec l'hydrure
de sodium, qui peut s'emballer au chauffage~; le DMF est en outre toxique
pour la reproduction. Mesures~: substituer le solvant (THF ou
2-méthyltétrahydrofurane) ou la base~; si on les garde, faire d'abord un
petit essai, à basse température, en ajoutant NaH par portions en suivant
la température, sous une hotte derrière un écran, avec un dégagement par
bulleur et un bain de refroidissement prêt~; mode opératoire écrit, une
seconde personne présente, et l'équipement de protection en dernier
recours.
\end{solution}

\begin{solution}{exo:b3:lab-techniques-3:11}
Contributions relatives~: 0.50, 0.50, $0.010/20.0 = 0.05$, $0.010/15.0 = 0.067$ et 0.10\,\%~; combinées,
$\sqrt{0.25 + 0.25 + 0.0025 + 0.0044 + 0.010} = 0.72$\,\%. Les deux intégrales dominent~: un meilleur rapport signal sur bruit
et des corrections soignées de ligne de base et de phase comptent plus que
la balance.
\end{solution}

\begin{solution}{exo:b3:lab-techniques-3:12}
$Q = (89.0 - 82.0)/(89.0 - 80.6) = 0.83 > 0.71$~: le test la signale. Avant de la rejeter, chercher dans le \omterm{def:b3:lab-techniques-3:notebook}{cahier
de laboratoire} une cause (une burette mal lue, une bulle d'air, un autre
lot de solution titrante)~; si l'on en trouve une, la rejeter en disant
pourquoi~; sinon, la garder et rapporter le test, ou refaire la mesure.
\end{solution}

\begin{solution}{pb:b3:lab-techniques-3:1}
\textbf{1.} $\ce{FeCl2 + 2NaC5H5 -> Fe(C5H5)2 + 2NaCl}$.

\textbf{2.} Le cyclopentadiénure de sodium est détruit par le dioxygène et
l'eau, et le chlorure de fer(II) s'oxyde en fer(III) à l'air humide.

\textbf{3.} $\qty{5.00}{g}/\qty{126.8}{g/mol} = \qty{0.0394}{mol}$~; $\times \qty{185.8}{g/mol} = \qty{7.33}{g}$~; rendement $5.86/7.33 = 80$\,\%.

\textbf{4.} $(0.10/1000)^3 = 10^{-12}$ (en théorie).

\textbf{5.} En chauffant le dimère, qui subit une rétro-Diels--Alder, et en
distillant le monomère volatil~; à température ambiante, le monomère se
redimérise par une réaction de Diels--Alder, si bien qu'on le garde au
froid et qu'on l'utilise aussitôt.

\textbf{6.} Chaque portion dégage du dihydrogène et de la chaleur~:
l'ajouter lentement permet de maîtriser le débit de gaz et la température.

\textbf{7.} \qty{186.0126}{u}~; erreur de $+1.4$ ppm.

\textbf{8.} C $120.0/185.8 = 64.58$\,\%, H $10.0/185.8 = 5.38$\,\%~; trouvé 64.41 et 5.47~: écarts de $-0.17$ et $+0.09$, à moins de 0.4.

\textbf{9.} Les dix hydrogènes sont équivalents par symétrie, et les cycles tournent rapidement.

\textbf{10.} Un seul signal, les dix carbones étant équivalents.

\textbf{11.} Une structure par diffraction des rayons X sur monocristal.

\textbf{12.} Non~: complexe à 18 électrons (\cref{ch:b3:organometallic-bonding}), il est stable à l'air, contrairement à ses
précurseurs.

\textbf{13.} Ses signaux ne chevauchent pas ceux du ferrocène, c'est un
solide pur, non volatil et non hygroscopique qui se pèse avec exactitude,
et il est inerte vis-à-vis de l'échantillon.

\textbf{14.} Une relaxation complète de chaque signal entre les acquisitions (un délai de plusieurs $T_1$),
une excitation uniforme, un rapport signal sur bruit suffisant, et des
corrections soignées de phase et de ligne de base.

\textbf{15.} Ferrocène $\ce{C10H10Fe}$ \qty{185.8}{g/mol}~; 1,3,5-triméthoxybenzène $\ce{C9H12O3}$ \qty{168.0}{g/mol}.

\textbf{16.} $P_x = 3.942 \times (3/10) \times (185.8/168.0) \times (15.0/20.0) \times 0.999 = 0.980$~: 98.0\,\%.

\textbf{17.} Relative~: $\sqrt{0.50^2 + 0.50^2 + 0.05^2 + 0.067^2 + 0.10^2} = 0.72$\,\%~; absolue~: $0.72\,\% \times 98.0 = 0.7$ point de pourcentage.

\textbf{18.} Oui, mais faiblement~: une impureté de 2\,\% de composition
voisine modifie C et H de moins que ce que l'analyse peut détecter~; la
qNMR est la mesure de pureté la plus informative.

\textbf{19.} L'hydrure de sodium (dihydrogène au contact de l'eau,
incendie), le dihydrogène dégagé, le THF et le cyclopentadiène (très
inflammables~; THF \omterm{def:b3:lab-techniques-3:peroxide}{peroxydable}), le craquage à chaud du dimère, le
chlorure de fer(II) (irritant) et le ferrocène (solide inflammable, nocif).

\textbf{20.} Hydrure de sodium~: pesé en dispersion dans l'huile sous gaz
inerte, ajouté par portions à une solution agitée et refroidie, le gaz
évacué par un bulleur, sans eau à proximité. THF~: sec et exempt de
peroxydes, issu d'une colonne, utilisé sous la hotte sous diazote.

\textbf{21.} Sous gaz inerte et avec refroidissement, par addition lente
d'isopropanol, puis d'éthanol, puis d'eau, en attendant chaque fois la fin
du dégagement gazeux.

\textbf{22.} Une vague réversible à un électron, $\ce{Fe(C5H5)2+}/\ce{Fe(C5H5)2}$, avec des pics distants d'environ 59 mV à \qty{25}{\celsius}
(\cref{ch:b3:electrode-kinetics})~; le couple est rapide, stable et presque
insensible au solvant, ce qui en fait une référence interne commode.

\textbf{23.} Le \omterm{def:b3:lab-techniques-3:notebook}{cahier de laboratoire}~: date, schéma, quantités et
origines, \omterm{def:b3:lab-techniques-3:assessment}{évaluation des risques}, ce qui a été fait et observé, rendements
et noms des fichiers de \omterm{def:b3:lab-techniques-3:notebook}{données brutes}. Les \omterm{def:b3:lab-techniques-3:literature}{informations
complémentaires}~: le mode opératoire complet et toutes les données de
caractérisation, avec des copies des spectres.

\textbf{24.} Pour que d'autres, et les auteurs des années plus tard,
puissent les rouvrir et les retraiter quel que soit leur logiciel~: des
données faciles à trouver, accessibles, interopérables et réutilisables.

\textbf{25.} La pureté qNMR de l'échantillon de ferrocène vaut $98.0 \pm 0.7$\,\% (incertitude-type).
\end{solution}
